\section{Four-Stage Validation Protocol}
\label{sec:protocol}

The protocol is ordered by the strength of evidence it yields: an environment
that merely runs (Stage~1) is not yet evidence that it teaches anything
(Stage~2), which is not yet evidence about \emph{what} is being optimized
(Stage~3), which is not yet evidence of causal capability gain (Stage~4). Each
stage is implemented in the released framework and produces one machine-readable
report.

\subsection{Stage 1: environment quality verification (no weight updates)}
\label{sec:stage1}

\paragraph{Goal.} Verify that generated environments are executable,
reproducible, and faithful to their source papers.

\begin{table}[t]
\centering
\small
\caption{Stage~1 metrics. ``Direction'' is stated explicitly because two of
these read in opposite senses: a higher shortcut \emph{success} rate is worse.
The symbol under each metric is its name in \texttt{validation/metrics.py}, so
any reported number can be checked against the implementation.}
\label{tab:stage1}
\begin{tabular}{@{}p{1.55in}p{3.5in}c@{}}
\toprule
Metric & Definition & Dir. \\
\midrule
Construction success rate\newline \texttt{build\_success\_rate}
  & Fraction of environments that build and run without error & $\uparrow$ \\
Baseline reproduction consistency\newline
  \texttt{baseline\_\allowbreak reproduction\_\allowbreak consistency}
  & Deviation from paper-reported results & $\uparrow$ \\
Expert approval rate\newline \texttt{expert\_approval\_rate}
  & Fraction of environments approved by human reviewers & $\uparrow$ \\
Shortcut success rate\footnotemark\newline \texttt{shortcut\_success\_rate}
  & Fraction of environments where the evaluator is deceived by a trivial
    strategy & $\downarrow$ \\
\bottomrule
\end{tabular}
\end{table}
\footnotetext{The draft nomenclature calls this ``shortcut detection rate''
while defining it as the fraction where the evaluator \emph{is deceived}---a
defect rate that decreases as quality improves. We keep the definition and
rename the quantity to match its direction; the implementation reports both
\texttt{shortcut\_success\_rate} ($\downarrow$) and
\texttt{shortcut\_detection\_rate} $= 1 - \text{success}$ ($\uparrow$) so either
convention can be read from the report without re-deriving it
(claim~\cid{C7}).}

Every Stage~1 metric returns \texttt{None}---not $0.0$---on empty input, so an
empty set can never be transcribed as a score of zero (claim~\cid{C5}).

\paragraph{Controls.} Expert-manually constructed environments
(DrugDiscoveryBench \cite{drugdiscoverybench}, Terminal-Bench-Science
\cite{tbscience}) serve as the gold standard; claims~\cid{P1}--\cid{P4} name the
control arm each rate is measured against.

\paragraph{Execution modes.} The runner consumes a per-environment results
manifest (offline scoring), can build and run each container directly when
Docker is present, or merges both (disk inventory for existence, external facts
for execution). When Docker is absent the affected environments are reported
\textsc{Unverified}; they are counted, never dropped and never counted as
successes.

\subsection{Stage 2: weight-free self-improvement (no weight updates)}
\label{sec:stage2}

\paragraph{Goal.} Verify that environments provide genuine improvement signals
across sequential attempts.

\paragraph{Protocol.} Following RSI-Forge's design, $N$ models perform $K$
sequential attempts per environment, inheriting code and notes but \emph{without
weight updates}. We measure the first-attempt improvement rate
(claim~\cid{P5}), the fraction of environments where some model exceeds the
reproduced baseline (claim~\cid{P6}), and the continuous improvement rate.
The scheduler is deterministic---environment-major, then round, then model---so a
run is reproducible, and it supports a dry-run mode that emits the full
execution plan without invoking any model.

\paragraph{Expected outcome.} Replication of RSI-Forge's finding of $\geq 84\%$
first-attempt improvement in life-science environments \cite{rsi-forge}.

\subsection{Stage 3: harness optimization (no weight updates)}
\label{sec:stage3}

\paragraph{Goal.} Verify that agents can improve through optimization of
external components---prompts, tools, memory, control flow---without weight
changes.

\paragraph{Protocol.} Evaluate on HarnessOpt-Bench-style tasks
\cite{harnessoptbench} measuring whether models can self-modify their execution
harness, and compare harness variants across OfficeQA \cite{officeqa}, GAIA
\cite{gaia}, and life-science \harbor\ tasks. The key finding to test is whether
model differences exceed harness differences, as HarnessOpt-Bench reports, or
whether the opposite holds in life-science domains where domain toolchains
(Bioconductor, samtools) may dominate (claim~\cid{P7}).

\paragraph{Implementation note.} The comparison decomposes observed spread into
between-model and between-harness components on a model $\times$ harness grid.
This decomposition is only interpretable on a \emph{complete rectangular} grid:
with missing cells the two components acquire different constant offsets and
``which axis dominates'' becomes an artifact of which cells happen to be filled.
We therefore require a complete grid before reporting the comparison; this is
recorded as a known limitation in Appendix~\ref{app:limits}.

\subsection{Stage 4: post-training with held-out benchmarks (weight updates)}
\label{sec:stage4}

\paragraph{Goal.} Produce causal evidence that generated environments improve
model capabilities on unseen tasks.

\paragraph{Training data.} Automatically generated \harbor\ environments from
papers published \textbf{before 2024}.

\paragraph{Held-out benchmark.} Environments generated from papers published
\textbf{2025--2026}, \emph{locked before training begins}. The lock records the
corpus digest, the split parameters, and the membership of every arm; a later
change to the corpus fails verification instead of silently drifting to match
whatever the training set ended up containing (claim~\cid{C6}). Supplement with
expert-manual tasks (DrugDiscoveryBench, Terminal-Bench-Science) and public
scientific benchmarks (BixBench \cite{bixbench}, LAB-Bench \cite{labbench},
ScienceAgentBench \cite{scienceagentbench}) after confirming no contamination.

\paragraph{Training methods.}
\begin{itemize}
  \item \textbf{RLVR}: for tasks with verifiable metrics (AUROC, F1, RMSE,
        peak-calling accuracy) \cite{rlvr,gvpo,entropyaware}.
  \item \textbf{RaR (Rubrics as Rewards)}: for composite or subjective tasks
        where no single standard answer exists \cite{rar}.
  \item \textbf{SFT / rejection sampling}: generate trajectories with a strong
        model and distill to the target model.
\end{itemize}

\begin{table}[t]
\centering
\caption{Stage~4 control arms. Every comparative claim names the arm it is
measured against (Appendix~\ref{app:claims}); a claim without a control is not
reportable in this manuscript.}
\label{tab:controls}
\begin{tabular}{@{}ll@{}}
\toprule
Control & Purpose \\
\midrule
No environment training & Baseline capability \\
Random environment training & Environment signal vs.\ generic training \\
Expert environment training & Upper bound from manual curation \\
Auto environment training & Effect of automation (\cid{P8}) \\
Auto + expert-reviewed training & Effect of expert review (\cid{P9}) \\
\bottomrule
\end{tabular}
\end{table}

\paragraph{Evaluation benchmarks.} Life-science agent benchmarks (BixBench,
LAB-Bench, ScienceAgentBench, DrugDiscoveryBench, Terminal-Bench-Science,
RSI-Exam, GAIA); self-built held-out \harbor\ environments from 2025--2026
papers; and real end-to-end workflows (single-cell analysis, variant
annotation, protein structure prediction).

\paragraph{Statistical rigor.} Multiple random seeds; confidence intervals;
compute-matched controls; ablations over environment count, domain diversity,
evaluator strictness, and expert review proportion; and a contamination check
confirming that benchmark tasks do not appear in training papers or environments
(claim~\cid{P10}).
